\section{功能：Long CoT 结构中各类键的塑形作用}
接下来我们分析不同键在 Long CoT 结构中的塑形功能。我们假设：LLM 在语义空间中寻找最优配置的过程，类似于蛋白沿折叠漏斗下滑到低能态（即解答）的过程。\vspace{-8pt}


\paragraph{\textbf{深层推理使主体结构致密。}} 随着推理推进，深层推理通过构造逻辑骨架来形成一级结构，类似共价键。图~\ref{fig:function}(a) 显示，深层推理会压缩核心逻辑结构：语义空间中最小包围球体积相较基线减少 22\%。这一步构建了答案的骨架，但尚未保证全局稳定或正确性。\vspace{-8pt}

\paragraph{\textbf{自我反思让全局逻辑致密且稳定。}} 在骨架延展之后，自我反思负责稳定结构。它类似氢键，通过跨越远距离的节点来检查一致性，而不是单纯增加推理步。图~\ref{fig:function}(b) 显示，反思会收紧“疏水核心”、抑制不一致分支，把系统体积从 35.2 降到 31.2，使结构走向稳定折叠态。\vspace{-8pt}

\paragraph{\textbf{自我探索扩展逻辑空间。}} 自我探索扩大了可行解空间，同时增加多样性、也可能引入不一致分支；它提高覆盖面，但会暂时降低稳定性。图~\ref{fig:function}(c) 表明，在学得 Long CoT 结构后，探索行为在语义空间中的覆盖体积从 23.95 扩大到 29.22。


\begin{TakeawayBox}{结论 5}
Long CoT 推理与蛋白折叠类似，可分为三个阶段：深层推理压实逻辑骨架，自我探索扩大搜索空间以避免局部最优，自我反思在语义空间中收敛到稳定、优化的解。
\end{TakeawayBox}
